\section{General measures with a model reference}
\label{sec:big-model-general-measures}

The projection produces a nonzero weighted measure below the prescribed datum. Maximize its mass among supersolutions with that same upper bound. If any part of the datum remains, weighted addition increases the mass and contradicts maximality. Density truncation then gives the equation for every finite non-pluripolar measure.

\begin{theorem}\label{thm:big-model-exponential}
Let $\mu$ be a nonzero finite non-pluripolar measure and let $\lambda>0$. There is a unique $u\in\mathcal E_\varphi$ such that
\[
 \mathrm{MA}_\theta(u)=\mathrm{e}^{\lambda u}\mu.
\]
\end{theorem}
\begin{proof}
We first take $\lambda=1$ and suppose that
\[
 \mu\leq A\sigma,\qquad A>0,\qquad
 \sigma=\mathrm{MA}_\omega(\psi),\quad
 \psi\in\PSH(X,\omega)\cap L^\infty(X).
\] The supersolution corollary gives $Q\in\mathcal E_\varphi$ with a finite nonzero measure
\[
 \nu_0=\mathrm{e}^{-Q}\mu_Q\leq\mu.
\]
Consider
\[
 \mathcal T=\{w\in\mathcal E_\varphi:
              \nu_0\leq \mathrm{e}^{-w}\mu_w\leq\mu\}.
\]
The family contains $Q$. Exponential order gives $w\leq Q$ for every $w\in\mathcal T$. Also
\[
 \ell\leq\mu_w(X)\leq \mathrm{e}^{\sup_Xw}\mu(X),
 \qquad
 \sup_Xw\geq\log\frac{\ell}{\mu(X)}.
\]
These bounds prevent collapse of any sequence in $\mathcal T$.

\Needspace{6\baselineskip}
Let
\[
 \alpha=\sup_{w\in\mathcal T}\int_X\mathrm{e}^{-w}\,\mu_w\leq\mu(X),
\]
and choose $w_j\in\mathcal T$ whose weighted masses tend to $\alpha$. Extract a genuine $L^1$ limit $w$. Since
\[
 \mu_{w_j}\leq \mathrm{e}^{\sup_XQ}\mu,
\]
the fixed-dominator theorem gives $w\in\mathcal E_\varphi$, convergence in capacity, and $\mu_{w_j}\rightharpoonup\mu_w$. Put $\nu_j=\mathrm{e}^{-w_j}\mu_{w_j}$. After another extraction, these finite measures converge weakly to a measure $\nu$, and
\[
 \nu_0\leq\nu\leq\mu,\qquad \nu(X)=\alpha.
\]

We claim $\mu_w=\mathrm{e}^w\nu$. \Cref{lem:fixed-npp-integrals} gives
\[
 \int_X|\mathrm{e}^{w_j}-\mathrm{e}^w|\,\mathrm{d}\mu\longrightarrow0.
\]
If $f$ is continuous, replacing $\mathrm{e}^{w_j}$ by $\mathrm{e}^w$ in $\int_Xf \mathrm{e}^{w_j}\,\mathrm{d}\nu_j$ therefore has error tending to zero. For the remaining integral, approximate the fixed bounded measurable function $f \mathrm{e}^w$ in $L^1(\mu)$ by continuous functions. Since $\nu_j,\nu\leq\mu$, the approximation error is uniform in $j$, while the integrals of each continuous approximation converge by weak convergence. Thus
\[
 \int_Xf\,\mu_{w_j}
 =\int_Xf \mathrm{e}^{w_j}\,\mathrm{d}\nu_j
 \longrightarrow\int_Xf \mathrm{e}^w\,\mathrm{d}\nu.
\]
This proves the claim. The measure $\nu$ ignores the poles of $w$, so $\mathrm{e}^{-w}\mu_w=\nu$. Hence $w\in\mathcal T$ attains the maximum $\alpha$.

If $\mu-\nu$ were nonzero, the supersolution corollary would supply $v\in\mathcal E_\varphi$ with
\[
 0\ne\nu_v=\mathrm{e}^{-v}\mu_v\leq\mu-\nu.
\]
In particular $\mu_v\leq \mathrm{e}^{\sup_Xv}\mu\leq A \mathrm{e}^{\sup_Xv}\sigma$, so the bounded-carrier hypothesis of the addition theorem is satisfied. That theorem gives $a\in\mathcal E_\varphi$ with
\[
 \mathrm{e}^{-a}\mu_a=\mathrm{e}^{-w}\mu_w+\mathrm{e}^{-v}\mu_v=\nu+\nu_v.
\]
The right side is between $\nu_0$ and $\mu$, but its mass exceeds $\alpha$. This contradicts maximality. Therefore $\nu=\mu$, and $w$ solves the equation.

\Needspace{8\baselineskip}
Now let $\mu$ be any nonzero finite non-pluripolar measure. The bounded-carrier decomposition writes
\[
 \mu=f\sigma,\qquad
 \sigma=\mathrm{MA}_\omega(\psi),\quad
 \psi\in\PSH(X,\omega)\cap L^\infty(X),\quad
 0\leq f\in L^1(\sigma).
\]
\Needspace{8\baselineskip}
Set $f_k=\min(f,k)$ and $\mu_k=f_k\sigma$. These measures are nonzero, and the first part gives $u_k\in\mathcal E_\varphi$ with $\mu_{u_k}=\mathrm{e}^{u_k}\mu_k$. Exponential order shows that $u_k$ decreases. The uniform bound
\[
 \sup_Xu_k\geq\log\frac{\ell}{\mu(X)}
\]
gives a genuine limit $u$. Moreover $\mu_{u_k}\leq \mathrm{e}^{\sup_Xu_1}\mu$. Fixed-dominator continuity gives $u\in\mathcal E_\varphi$ and weak convergence of these measures. On the right,
$\mathrm{e}^{u_k}f_k\to \mathrm{e}^u f$ in $L^1(\sigma)$ by dominated convergence, with dominator $\mathrm{e}^{\sup_Xu_1}f$. Thus $\mu_u=\mathrm{e}^u\mu$.

For general $\lambda>0$, apply the exponent-one result in background $\lambda\theta$, with model $\lambda\varphi$ and datum $\lambda^n\mu$, and divide its potential by $\lambda$. Scaling of the class and the Monge--Amp\`ere measure gives the displayed equation. Finally, if $u,v\in\mathcal E_\varphi$ solve it, then $\mu_u=\mathrm{e}^{\lambda(u-v)}\mu_v$. The strict domination theorem with constant shifts gives $u=v$.
\end{proof}

\Needspace{15\baselineskip}
\begin{theorem}\label{thm:big-model-normalized}
For every nonzero finite non-pluripolar measure $\mu$, there are $u\in\mathcal E_\varphi$ and $c>0$ such that
\[
 \sup_Xu=0,\qquad \mathrm{MA}_\theta(u)=c\mu.
\]
The scalar $c$ is unique and satisfies
\[
 \frac{\ell}{\mu(X)}\leq c\leq\frac{M}{\mu(X)}.
\]
\end{theorem}
\begin{proof}
For each integer $j\geq1$, solve
\[
 \mu_{u_j}=\mathrm{e}^{u_j/j}\mu,\qquad u_j\in\mathcal E_\varphi,
\]
using the preceding theorem. Put $S=\sup_Xu_1$ and $w_j=u_1+(j-1)S$. Then $w_j/j\geq u_1$, and hence
\[
 \mu_{w_j}=\mathrm{e}^{u_1}\mu\leq \mathrm{e}^{w_j/j}\mu.
\]
Exponential order at exponent $1/j$ gives $u_j\leq w_j$. Define
\[
 v_j=u_j-\sup_Xu_j,\qquad
 c_j=\mathrm{e}^{(\sup_Xu_j)/j}.
\]
The upper bound just proved and the positive lower mass imply
\[
 \frac{\ell}{\mu(X)}\leq c_j\leq \mathrm{e}^S,
 \qquad \mu_{v_j}=c_j \mathrm{e}^{v_j/j}\mu\leq \mathrm{e}^S\mu.
\]
The envelope identity gives $v_j\leq\varphi$. \Cref{thm:normalized-compactness} supplies a subsequence $j_k$ for which $v_{j_k}\to u$ in $L^1$, with $u$ genuine and $\sup_Xu=0$, and $c_{j_k}\to c>0$. Fixed-dominator continuity gives $u\in\mathcal E_\varphi$ and $\mu_{v_{j_k}}\rightharpoonup\mu_u$.

\Cref{lem:normalized-exponential-npp} gives $\mathrm{e}^{v_j/j}\to1$ in $L^1(\mu)$, also for this non-semipositive background. Consequently $c_{j_k}\mathrm{e}^{v_{j_k}/j_k}\mu\to c\mu$ in total variation. The retained reciprocal exponents are $1/j_k$, so the limit equation is $\mu_u=c\mu$. Integrating it gives the stated interval for $c$. The scalar comparison in \cref{cor:big-relative-orders} gives uniqueness of $c$.
\end{proof}

For $L^p$ data, \cref{sec:big-model-lp} will prove the stronger conclusion $u\simeq\varphi$.

\subsection{Stability under changes of reference}
\label{subsec:big-equation-reference-stability}

Convergence of the models in $d_\theta$ controls the relative class of a limit. For equations, convergence of the data supplies the fixed dominating measures needed after extraction. The positive exponent then identifies the potential; at exponent zero, scalar uniqueness identifies the limiting measure.

\begin{theorem}\label{thm:big-varying-equation-stability}
Let $\varphi_j,\varphi\in\mathcal M_\delta(\theta)$ for some $\delta>0$, with $d_\theta(\varphi_j,\varphi)\to0$. Let $\nu$ be a finite non-pluripolar measure and let $f_j,f$ be nonnegative nonzero functions in $L^1(\nu)$ with $f_j\to f$ in $L^1(\nu)$. Put $\mu_j=f_j\nu$ and $\mu=f\nu$.
\begin{enumerate}
\item For a fixed $\lambda>0$, the solutions
\[
 u_j\in\mathcal E_{\varphi_j},\qquad
 \mathrm{MA}_\theta(u_j)=\mathrm{e}^{\lambda u_j}\mu_j
\]
converge in $L^1(X)$ and in capacity to the solution $u\in\mathcal E_\varphi$ with datum $\mu$. Their Monge--Amp\`ere measures converge to $\mathrm{MA}_\theta(u)$ in total variation.
\item Let $u_j\in\mathcal E_{\varphi_j}$ satisfy $\sup_Xu_j=0$ and $\mathrm{MA}_\theta(u_j)=c_j\mu_j$. Then $c_j\to c$, where $c$ is the scalar for reference $\varphi$ and datum $\mu$, and
\[
 \mathrm{MA}_\theta(u_j)\longrightarrow c\mu
 \quad\text{in total variation}.
\]
The potentials form an $L^1$ relatively compact family. Every $L^1$ cluster limit $u$ lies in $\mathcal E_\varphi$, has supremum zero and solves $\mathrm{MA}_\theta(u)=c\mu$. Each subsequence converging in $L^1$ converges to that limit in capacity.
\end{enumerate}
\end{theorem}
\begin{proof}
Write $M_0=\overline{\operatorname{vol}}_\theta(V_\theta)$. Every solution above has mass between $\delta$ and $M_0$. In the first case let $s_j=\sup_Xu_j$. Then
\[
 s_j\geq\lambda^{-1}\log\frac{\delta}{\mu_j(X)},
\]
which is a uniform lower bound. To bound $s_j$ above, suppose after extraction that $s_j\to+\infty$. Normalized compactness gives $v_j=u_j-s_j\to v$ in $L^1$ with $\sup_Xv=0$. The fixed-measure lemma and $f_j\to f$ in $L^1(\nu)$ imply
\[
 \int_X\mathrm{e}^{\lambda v_j}f_j\,\mathrm{d}\nu
 \longrightarrow\int_X\mathrm{e}^{\lambda v}f\,\mathrm{d}\nu>0.
\]
The limit is positive since $f\nu$ is nonzero and ignores the poles of $v$. But the left side equals $\mathrm{e}^{-\lambda s_j}\mu_{u_j}(X)\leq M_0\mathrm{e}^{-\lambda s_j}$, a contradiction. Thus $s_j$ is bounded and the $u_j$ form an $L^1$ relatively compact family of genuine potentials.

Start with any subsequence and extract $u_j\to u$ in $L^1$. Extract again so that $\sum_j\|f_j-f\|_{L^1(\nu)}<\infty$. The function
\[
 g=f+\sum_j|f_j-f|\in L^1(\nu)
\]
dominates all retained $f_j$. If $C$ is their common upper bound for $u_j$, then $\mu_{u_j}\leq\mathrm{e}^{\lambda C}g\nu$. By \cref{thm:big-varying-model-closure}, $u\in\mathcal E_\varphi$, the potentials converge in capacity and the measures converge weakly to $\mu_u$. On the other hand,
\[
 \mathrm{e}^{\lambda u_j}f_j\longrightarrow
 \mathrm{e}^{\lambda u}f\quad\text{in }L^1(\nu).
\]
Indeed, the error from $f_j-f$ is at most $\mathrm{e}^{\lambda C}\|f_j-f\|_1$, and \cref{lem:fixed-npp-integrals} for the fixed measure $f\nu$ controls the other term. Hence $\mu_u=\mathrm{e}^{\lambda u}\mu$, with convergence in total variation. Uniqueness in \cref{thm:big-model-exponential} identifies every cluster limit. The subsequence criterion gives all three convergences for the original sequence.

For the second case, integrating the equation gives
\[
 \frac{\delta}{\mu_j(X)}\leq c_j\leq\frac{M_0}{\mu_j(X)}.
\]
Since $\mu_j(X)\to\mu(X)>0$, these scalars stay in a compact subinterval of $(0,+\infty)$. Normalized compactness gives, after extraction, $u_j\to u$ in $L^1$, $\sup_Xu=0$, and $c_j\to c'>0$. Extract a summable density error as above. Then $\mu_{u_j}\leq Cg\nu$ for a fixed $C$, so \cref{thm:big-varying-model-closure} gives $u\in\mathcal E_\varphi$, capacity convergence and weak convergence to $\mu_u$. Since $c_jf_j\to c'f$ in $L^1(\nu)$, the limit equation is $\mu_u=c'\mu$. Scalar uniqueness in \cref{thm:big-model-normalized} gives $c'=c$. Thus the whole scalar sequence converges, and so does $c_j\mu_j$ in total variation. Applying the same argument within any $L^1$ convergent subsequence proves its capacity convergence and the asserted cluster equation.
\end{proof}
